%========================================================================================
%  CAPITOLO 3 — IL MODELLO DI RUMORE DI UNA PULSAR
%  Obiettivo: descrivere tutte le componenti del modello usato nelle analisi.
%  Lunghezza indicativa: 5-6 pagine.
%  Fonti: Taylor (2021) cap. 7, Shaifullah (2025) sez. 3, EPTA DR2 paper II (EPTA2023_noise),
%         Lentati et al. (2016), tesi di Irene sez. 2.1.
%========================================================================================

\chapter{Il modello di rumore di una pulsar} \label{chpt:rumore}
\sloppy

% Intro: i residui = segnale + rumore; il rumore si divide in bianco e rosso; il rosso si
% divide in acromatico e cromatico. Tabella riassuntiva dei 6 parametri usati in tesi.

\section{Rumore bianco} \label{sec:rumore_bianco}
% - Incertezza radiometrica sui TOA, jitter.
% - Parametri per backend: EFAC (riscala le barre d'errore), EQUAD/TNEQUAD (aggiunge in
%   quadratura), ECORR (correlato tra TOA della stessa osservazione).
% - Formula: sigma^2 = EFAC^2 * (sigma_TOA^2 + EQUAD^2).
% - Nella tesi: fissati ai valori del noisedict (marginalizzati separatamente).

\section{Rumore rosso acromatico} \label{sec:rumore_rosso}
% - Spin noise: irregolarità della rotazione. Non dipende dalla frequenza radio.
% - Legge di potenza P(f) = A^2/(12 pi^2) (f/fyr)^(-gamma) fyr^(-3)
%   -> stessa forma del GWB (Eq. psd_gwb)! Parametri: red_noise_log10_A, red_noise_gamma.

\section{Rumore cromatico: dispersione e scattering} \label{sec:rumore_cromatico}
% - Variazioni di DM nel tempo (moto relativo pulsar-Terra-ISM): ritardo prop. a nu^-2.
%   Parametri: dm_gp_log10_A, dm_gp_gamma.
% - Scattering nel mezzo interstellare: ritardo prop. a nu^-alpha, alpha = 4 (spettro di
%   Kolmogorov darebbe 4.4). Parametri: sv_gp_log10_A, sv_gp_gamma (+ sv_gp_alpha se libero).
% - Perché serve una buona copertura in frequenza radio (Ferranti2025).

\section{Processi gaussiani e base di Fourier} \label{sec:GP}
% - Ogni processo rosso scritto come somma di seni e coseni alle frequenze k/T,
%   k = 1..N, con coefficienti gaussiani di varianza data dalla PSD.
% - Per i processi cromatici la base è moltiplicata per (nu/nu_ref)^(-alpha).
% - Numero di componenti di Fourier per J1640+2224 (dal noisedict di DR2low).
% - Figura: PSD delle tre componenti + white noise per J1640 (stile Irene fig. 2.7).

\section{Riepilogo del modello} \label{sec:riepilogo_modello}
% Tabella: componente | dipendenza da nu | parametri | prior.
